\documentclass[11pt]{article}
\usepackage{mathblatt}
\makeatletter
\setlength{\mbpfbe}{0mm}% keine Punkte (Bauregel 6.4)
\makeatother
% Rechenraum als Karo (Bauregel 6.9): n Rechenschritte = n mal 10 mm
\newcommand{\karo}[1]{\par\smallskip\noindent\begin{tikzpicture}
  \clip (0,0) rectangle ({\linewidth-1mm},{#1*10mm});
  \draw[black!16,line width=0.3pt,step=5mm] (0,0) grid ({\linewidth},{#1*10mm});
  \end{tikzpicture}\par}
\newcommand{\linie}[1]{\,{\color{black!45}\rule[-1pt]{#1}{0.4pt}}\,}
\newcommand{\pkt}{\textperiodcentered{}}
\begin{document}

\pfheftstil
\pflaufkopf[P10 \pkt{} mehr]{Quadratische Gleichung lösen \pkt{} Lösungen}{PA40-3}\pfseitevon
\pfheftkopf{Quadratische Gleichung lösen}{Lösungen \pkt{} mehr}
\begin{pfloesungb}{}
\lzb{1.}{}{x$_1$ = 2; x$_2$ = $-$5}{x$_1$ = 2; x$_2$ = $-$5}
\end{pfloesungb}
\begin{pfloesungb}{}
\lzb{2.}{}{(x + 4)² $-$ 3 = x² + 8x + 13; x $\approx$ $-$2,27; x$_2$ $\approx$ $-$5,73}{(x + 4)² $-$ 3 = x² + 8x + 13; x = $-$4 ± √3, also x$_1$ $\approx$ $-$2,27; x$_2$ $\approx$ $-$5,73}
\end{pfloesungb}
\end{document}